\documentclass[main.tex]{subfiles}
\begin{document}

\section{Exercise 1 – Operating Region of the NMOS Transistor}

\subsection*{Given}
\[
\begin{aligned}
\mu_n C_{ox} &= 110~\mu\text{A/V}^2, &
\frac{W}{L} &= 4, &
V_T &= 1.0~\text{V},\\
\lambda &= 0.1~\text{V}^{-1}, &
R_O &= 2.2~\text{k}\Omega, &
I_{in} &= 1.09~\text{mA}, &
V_{EE} &= -4.96~\text{V}.
\end{aligned}
\]

\subsection*{Method}
For an NMOS transistor in saturation:
\begin{equation}
I_{in} = \tfrac{1}{2}\mu_n C_{ox}\tfrac{W}{L}(V_{GS}-V_T)^2(1+\lambda V_{DS})
\label{eq:1.1}
\end{equation}
and the saturation condition:
\begin{equation}
V_{DS} \ge V_{GS}-V_T.
\label{eq:1.2}
\end{equation}
Neglecting channel-length modulation:
\begin{equation}
V_{\text{eff}} = V_{GS}-V_T = \sqrt{\tfrac{2I_{in}}{\mu_n C_{ox}(W/L)}}.
\label{eq:1.3}
\end{equation}

\subsection*{Calculation}
\[
V_{\text{eff}} = \sqrt{\tfrac{2(1.09\times10^{-3})}{110\times10^{-6}\times4}}
=2.23~\text{V},\qquad
V_{GS}=V_T+V_{\text{eff}}=1+2.23=3.23~\text{V}.
\]
With \(V_D\approx0\) and \(V_S=V_{EE}=-4.96~\text{V}\):
\[
V_{DS}=V_D-V_S=0-(-4.96)=4.96~\text{V}>V_{\text{eff}}=2.23~\text{V}.
\]

\subsection*{Result}
\[
\boxed{\text{The NMOS transistor operates in the saturation (active) region.}}
\]

\subsection*{Note}
\(C_{gs}\), \(C_{gd}\), and \(C_{db}\) affect only AC behaviour
and are omitted in this DC analysis.

\end{document}
